%%%%%%%%%%%%%%%%%%%%%%%%%%%%%%%%%%%%%%%%%%%%%%%%%%%%%%%
% % Lines starting with % are comments, which are ignored.
% % This is a handy way of indicating the date and version of
% % your document, to wit:
% %
% % LaTeX sample file
% % Modified March, 2002
% %
%%%%%%%%%%%%%%%%%%%%%%%%%%%%%%%%%%%%%%%%%%%%%%%%%%%%%%%
\DocumentMetadata{
  % PDF standards: ONE key whose value is a comma-separated list --
  % UA-2 (accessibility) plus A-4f (archival; the archival companion of
  % UA-2). Do NOT use separate pdfstandard=... keys: with ordinary
  % key-value parsing only the last one takes effect.
  pdfstandard={UA-2,A-4f},
  lang=en-US,
  tagging=on,
  tagging-setup={
    math/setup=mathml-SE
  }
}

\documentclass{ltx-talk}
% lecture-macros.tex -- shared package setup for lecture-template.tex
% (ltx-talk slides / handouts).
%
% NOTE ON PATHS: every template in this repository is stored one level up
% from where it is compiled (see README.md). All relative \input paths
% below are resolved against the folder the document is compiled in,
% which sits one directory above this file.

% --- Tables, bookmarks, strikethrough ------------------------------------
% booktabs: professional table rules (\toprule, \midrule, \bottomrule)
% bookmark: PDF outline entries (PDF bookmarks)
% ulem[normalem]: \sout strikethrough that survives math mode
\usepackage{booktabs}
\usepackage{bookmark}
\usepackage[normalem]{ulem}
\usepackage{multicol}

% --- Math ------------------------------------------------------------
% amsmath: equation environments and numbering
% mathrsfs: \mathfrak  |  stmaryrd: extra symbols
\usepackage{amsmath}
\usepackage{mathrsfs}
\usepackage{stmaryrd}

% --- Layout -------------------------------------------------------------
% float: [H] float placement
% fontspec: Unicode/OpenType fonts (compile with XeLaTeX or LuaLaTeX)
\usepackage{float}
\usepackage{fontspec}

% --- Colors ---------------------------------------------------------------
% dvipsnames: extended named color palette (the Accessible* colors below
% are built on it);
% table: \rowcolors / \cellcolor for shading tables.
% NOTE: xcolor may only be loaded ONCE per document -- a second
% \usepackage{xcolor} with different options is a fatal "Option clash"
% error, so ALL of its options are listed here.
\usepackage[dvipsnames,table]{xcolor}

% --- Multimedia -----------------------------------------------------------
% animate: embeds audio/video (and legacy Flash) media in a frame.
\usepackage{animate}

% --- Bibliography ----------------------------------------------------------
% biblatex (biber backend): citations for lecture notes.
\usepackage[
  backend=biber,
  style=alphabetic,
  url=true,
  doi=true,
  isbn=true,
  backref=true
]{biblatex}

% --- Colored boxes ----------------------------------------------------------
% talkthemetcolorbox: the colored-box theme for ltx-talk slides
% (shadow/rounded/shaded blocks, showtitle, titlepage styles).
\usepackage[
  shadow,
  rounded,
  shaded,
  blocks,
  showtitle,
  titlepage
]{talkthemetcolorbox}

% --- "theindex" helper environment ------------------------------------------
% Provides a lightweight theindex environment (an itemize) and the
% \@idxitem/\subitem/\subsubitem/\indexspace macros that LaTeX's standard
% theindex machinery expects, so hand-written index pages indent correctly.
\makeatletter
\newenvironment{theindex}
  {\begin{itemize}}
  {\end{itemize}}
\providecommand\@idxitem{\par\hangindent 40\p@}
\providecommand\subitem{\@idxitem \hspace*{20\p@}}
\providecommand\subsubitem{\@idxitem \hspace*{30\p@}}
\providecommand\indexspace{\par \vskip 10\p@ \@plus5\p@ \@minus3\p@\relax}
\makeatother

% --- Index -----------------------------------------------------------------
% imakeidx: builds the printed index.
% The redefinition below appends "|hyperpage" to every index entry so that,
% in the PDF, each index entry is a link to its page.
\usepackage{imakeidx}
\makeatletter
\let\imki@orig@wrindexentrysplit\imki@wrindexentrysplit
\renewcommand\imki@wrindexentrysplit[3]{%
  \imki@orig@wrindexentrysplit{#1}{#2|hyperpage}{#3}%
}

\makeatletter
% --- Section titles: capture + one bookmark per logical section -------------
% \section is redefined so that:
%   1. the section title is captured into the plain-TeX register
%      \mytalksectitletoks (reprinted verbatim on the auto divider frame,
%      see the cmd/section/after hook below -- plain-TeX storage is
%      catcode-safe, unlike expanding the title in a hook), and
%   2. exactly ONE depth-1 PDF bookmark is created per \section call
%      (hyperref's own bookmark is suppressed via \Hy@writebookmark).
\DeclareCommandCopy{\OriginalSectionForBookmarks}{\section}
\newtoks\mytalksectitletoks

\RenewDocumentCommand{\section}{s o m}{%
  \begingroup
    \def\Hy@writebookmark##1##2##3##4##5{}%
    \global\mytalksectitletoks={#3}%
    \IfBooleanTF{#1}
      {\phantomsection\IfNoValueTF{#2}{\OriginalSectionForBookmarks*{#3}}{\OriginalSectionForBookmarks*[#2]{#3}}}%
      {\IfNoValueTF{#2}{\OriginalSectionForBookmarks{#3}}{\OriginalSectionForBookmarks[#2]{#3}}}%
  \endgroup
  \IfBooleanTF{#1}
    {\IfNoValueTF{#2}{\pdfbookmark[1]{#3}{section-star.\thepage.\number\inputlineno}}{\pdfbookmark[1]{#2}{section-star.\thepage.\number\inputlineno}}}%
    {\IfNoValueTF{#2}{\pdfbookmark[1]{Section \thesection: #3}{section.\thesection}}{\pdfbookmark[1]{Section \thesection: #2}{section.\thesection}}}%
}
\makeatother

% --- One PDF bookmark per logical frame (NOT per overlay) ---------------------
% ltx-talk writes one PDF page per overlay, so a naive frametitle hook would
% register a new bookmark for every overlay of the same frame. The
% \g_slideutils_frame_bookmark_int register remembers the last frame that
% got a bookmark, so subsequent overlays of the same frame are skipped.
% This is also what keeps BeamerPresenter's cumulative annotations working:
% one logical frame must stay one logical unit in the PDF outline.
\ExplSyntaxOn
\int_new:N \g_slideutils_frame_bookmark_int

\AddToHook{cmd/frametitle/after}{%
  \tl_if_empty:NF \g__talk_frame_title_tl
    {
      \ifnum\value{frame}=\g_slideutils_frame_bookmark_int\relax
      \else
        \int_gset:Nn \g_slideutils_frame_bookmark_int { \value{frame} }%
        \pdfbookmark[2]{\tl_use:N \g__talk_frame_title_tl}
          {frame.\int_use:N \g_slideutils_frame_bookmark_int}
      \fi
    }
}
\ExplSyntaxOff
\makeatother

% --- Printed index + theme for the colored boxes -------------------------------
% \makeindex[...]: two-column index titled "Index", listed in the TOC.
% The two EditInstance calls style the standard header/footer instances
% of talkthemetcolorbox (cyan background; footer order title/author/
% date/framenumber).
\makeindex[columns=2, title=Index, intoc]\EditInstance{header}{std}{
  background-color = cyan,
  color            = black,
  height           = 0.75cm
}
\EditInstance{footer}{std}{
  background-color = cyan,
  element-order = {title, author, date, framenumber},
}

% --- Auto section divider frame ---------------------------------------------------
% After each \section, insert a centered, large, structure-colored frame
% that prints the section title (from \mytalksectitletoks).
\AddToHook{cmd/section/after}{%
  \begin{frame}
    \centering
    \color{structure}
    \LARGE
    \the\mytalksectitletoks
  \end{frame}%
}

% --- Tagging sockets (tagged PDF structure) ---------------------------------------
% ltx-talk exposes tagging "sockets" where the PDF structure can attach
% section titles; register one and leave it unassigned (none) so the
% class's own tagging keeps control.
\ExplSyntaxOn
\AtBeginDocument{
  \NewTaggingSocketPlug { talk / sec / title } { none } { }
  \AssignTaggingSocketPlug { talk / sec / title } { none }
}
\ExplSyntaxOff

% --- Accessibility roles ----------------------------------------------------------------
% Map \frametitle to the H1 heading role in the tagged PDF, so screen
% readers treat slide titles as headings.
\tagpdfsetup{role / new-tag = frametitle / H1}

\makeatletter
% --- Page label = frame number -------------------------------------------------------
% Redefine what feeds the PDF page label to track the frame counter,
% not the physical page counter. Replace \theframenumber below with
% whatever internal macro/property ltx-talk actually uses to produce
% the "framenumber" footer element — confirm via ltx-talk-frame.dtx.
\renewcommand\thepage{\arabic{frame}}
\makeatother

\addbibresource{refs.bib}

% --- "Problem Time!" slides ------------------------------------------------------------------
% \ProblemSlide{problem}{solution} emits two frames: the problem, then
% its solution after a \pause, each numbered by the problem counter.
\newcounter{problemcounter}
\setcounter{problemcounter}{0}
\newcommand{\ProblemSlide}[2]{
  \stepcounter{problemcounter}
  \begin{frame}
    \frametitle{Problem Time!}
    Problem \theproblemcounter: #1
  \end{frame}

  \begin{frame}
    \frametitle{Problem Time!}
    Solution to Problem \theproblemcounter: \pause

    #2
  \end{frame}
}
% \highlight[<colour>]{<stuff>}
% from https://tex.stackexchange.com/questions/33401/a-version-of-colorbox-that-works-inside-math-environments
\newcommand{\highlight}[2][yellow]{\mathchoice%
  {\colorbox{#1}{$\displaystyle#2$}}%
  {\colorbox{#1}{$\textstyle#2$}}%
  {\colorbox{#1}{$\scriptstyle#2$}}%
  {\colorbox{#1}{$\scriptscriptstyle#2$}}}%

% --- Fonts ------------------------------------------------------------------------------
% Lato (text) + LeteSansMath (math): open, accessible typefaces.
\setmainfont{Lato}[
  Ligatures=TeX,
  Numbers={Lining,Monospaced}
]
\setsansfont{Lato}[
  Ligatures=TeX,
  Numbers={Lining,Monospaced}
]
\setmathfont{LeteSansMath.otf}\usepackage{emoji}

% --- Accessible colors ---------------------------------------------------------------------
% WCAG-contrast-safe variants of red/green for text on white:
% red!80!black and green!40!black both clear the 4.5:1 requirement.
% AltYellow: soft background for highlight boxes.
\colorlet{AccessibleRed}{red!80!black}
\colorlet{AccessibleGreen}{green!40!black}
\definecolor{AltYellow}{RGB}{255,245,170}

% --- Links ---------------------------------------------------------------------------------
% hyperref + link colors: every link uses AccessibleGreen (7.2:1 on
% white), so internal, citation, and URL links all pass WCAG 1.4.3.
\usepackage{hyperref}
\hypersetup{
  colorlinks=true,
  linkcolor=AccessibleGreen,
  urlcolor=AccessibleGreen,
  citecolor=AccessibleGreen,
}

% --- Shared course macros ---------------------------------------------------------------------
% Root macros.tex: math shorthands used across all course documents
% (\RR, \NN, \ZZ, \QQ, \grad, \curl, \tr, \der, \pder, ...).
\newcommand{\RR}{\mathbb{R}}
\newcommand{\NN}{\mathbb{N}}
\newcommand{\ZZ}{\mathbb{Z}}
\newcommand{\QQ}{\mathbb{Q}}
\newcommand{\vv}{\vec{v}}
\newcommand{\ww}{\vec{w}}
\newcommand{\xx}{\vec{x}}
\newcommand{\yy}{\vec{y}}
\newcommand{\zz}{\vec{z}}
\newcommand{\uu}{\vec{u}}
\newcommand{\basis}[1]{\vec{e_{#1}}}
\newcommand{\id}{\mathrm{id}}
\newcommand{\CC}{\mathbb{C}}
\renewcommand{\div}{\operatorname{div}}
\newcommand{\grad}{\operatorname{grad}}
\newcommand{\curl}{\operatorname{curl}}
\newcommand{\tr}{\operatorname{tr}}
\newcommand{\der}[2]{\dfrac{d #1}{d #2}}
\newcommand{\pder}[3][]{\dfrac{\partial^{#1} #2}{\partial #3}}


% % Title and author(s)
%%%%%%%%%%%%%%%%%%%%%%%%%%%%%%%%%%%%%%%%%%%%%%%%%%%%%%%
\title{Lecture 6 - The Derivative as a Function}
\author{Garrett Credi}
\date{09/24/2026}
%%%%%%%%%%%%%%%%%%%%%%%%%%%%%%%%%%%%%%%%%%%%%%%%%%%%%%%
% %
% % The next command allows your in import encapsulated
% % postscript files, .epsf or .eps files, which
% % contain vector graphic image data.
% %
%%%%%%%%%%%%%%%%%%%%%%%%%%%%%%%%%%%%%%%%%%%%%%%%%%%%%%%
% %   The first thanks indicates your affiliation
% %
% %  Just the name here.
% %
% % Your mailing address goes at the end.
% %
% % \thanks is also how you indicate grant support
% %
%%%%%%%%%%%%%%%%%%%%%%%%%%%%%%%%%%%%%%%%%%%%%%%%%%%%%%%

%%%%%%%%%%%%%%%%%%%%%%%%%%%%%%%%%%%%%%%%%%%%%%%%%%%%%%%
% $ The next command ensures the margins are 1 inch on
% % all sides.
% %
%%%%%%%%%%%%%%%%%%%%%%%%%%%%%%%%%%%%%%%%%%%%%%%%%%%%%%%
\begin{document}
\begin{frame}
  \maketitle
\end{frame}

\begin{frame}
  \frametitle{Quiz Time!}

  Please get into organized rows so that you have a ``bubble" of space around you
  away from other people.

  Reminder:
  \begin{enumerate}
    \item Do not start the quiz until I announce you can start.
    \item Do not read off other students' work to cheat
    \item Place all class material (notes, worksheets, etc) and devices (phones, tablets, headphones, etc) away
    \item Don't cheat and or be shady... obviously...
  \end{enumerate}
\end{frame}

\begin{frame}
  \frametitle{Last Time}

  Last lecture we: \pause
  \begin{enumerate}
    \item Introduced the idea of linearly approximating functions as an improvement to the constant-value approximation $f(x) \approx f(a)$ for $x \approx a$
      of continuous functions. \pause
    \item Defined the \textit{derivative} of a function $f(x)$ at a point $x = a$ as 
      \begin{equation} \pause
        f'(a) = \lim\limits_{x \rightarrow a} \dfrac{f(x) - f(a)}{x - a} = \lim\limits_{h \rightarrow 0} \dfrac{f(a + h) - f(a)}{h}
      \end{equation} \pause
    \item Saw that this quantity was also: \pause
      \begin{enumerate}
        \item The slope of the best linear approximation 
          \begin{equation}
            \mathscr{L}_{x = a}(x) = f(a) + f'(a)(x - a)
          \end{equation}
          to $f(x)$ at $x = a$. \pause
        \item The slope of the tangent line to the graph of $y = f(x)$ at $(a, f(a))$. \pause
        \item The instantaneous velocity of a particle whose motion is described by $f(x)$ at $x = a$.
      \end{enumerate}
  \end{enumerate}
\end{frame}

\begin{frame}
  \frametitle{Last Time Cont.}
  We also were able to compute some derivatives using strategies from our limit section of the course \pause e.g. for $f(x) = x^2$: \pause
  \begin{columns}
    \begin{column}{0.32\textwidth}
      \begin{align}
        f'(0) &= \pause \lim\limits_{x \rightarrow 0} \dfrac{f(x) - f(0)}{x - 0} \pause \\
              &= \lim\limits_{x \rightarrow 0} \dfrac{x^2}{x} \pause \\
              &= \lim\limits_{x \rightarrow 0} x \pause = 0
      \end{align}
    \end{column}

    \begin{column}{0.32\textwidth}
      \begin{align}
        f'(1) &= \pause \lim\limits_{x \rightarrow 1} \dfrac{f(x) - f(1)}{x - 1} \pause \\
              &= \lim\limits_{x \rightarrow 1} \dfrac{x^2 - 1}{x - 1} \pause \\
              &= \lim\limits_{x \rightarrow 1} x + 1 \pause = 1 + 1 = 2
      \end{align}
    \end{column}

    \begin{column}{0.32\textwidth}
      \begin{align}
        f'(2) &= \pause \lim\limits_{x \rightarrow 2} \dfrac{f(x) - f(2)}{x - 2} \pause \\
              &= \lim\limits_{x \rightarrow 2} \dfrac{x^2 - 4}{x - 2} \pause \\
      \end{align}
      \pause
    \end{column}
  \end{columns}

  Wait... these all seem pretty similar!
\end{frame}

\section{The Derivative as a Function}

\begin{frame}
  \frametitle{The Derivative Function I}
  What kind of mathematical object do we use
  when we want to compactly describe
  the process of taking some input value $a$ \pause
  and computing an output quantity $f'(a)$? \pause

  A function! \pause

  Let's define \textit{a derivative function} to see if we can simplify and speed up computations: \pause
  \begin{block}{The Derivative Function}\index{derivative!as a function}
    For a function $f(x)$ \pause 
    we define the \textit{derivative} (function), denoted either as: \pause
    \begin{center}
      $f'(x)$ or \pause $\dfrac{df}{dx}$ or \pause $\dfrac{d}{dx}f(x)$,
    \end{center}\pause
    to be the function so that: \pause
    \begin{equation}
      f'(x) = \lim_{h \rightarrow 0} \dfrac{f(x + h) - f(x)}{h} = \lim_{u \rightarrow x} \dfrac{f(u) - f(x)}{u - x}
    \end{equation} \pause
  \end{block}

  Is this different from what we were doing before? 
\end{frame}

\begin{frame}
  \frametitle{The Derivative Function II}
  Yes! We can now compute \textit{every} value of a derivative at once! \pause E.g. for $f(x) = x^2$ again

  \begin{align}
    f'(x) &= \pause \lim\limits_{h \rightarrow 0} \dfrac{f(x + h) - f(x)}{h} \pause \\
          &= \lim\limits_{h \rightarrow 0} \dfrac{(x+h)^2 - x^2}{h} \pause \\
          &= \lim\limits_{h \rightarrow 0} \dfrac{x^2 + 2xh + h^2 - x^2}{h} \pause \\
          &= \lim\limits_{h \rightarrow 0} \dfrac{2xh  + h^2}{h} \pause \\
          &= \lim\limits_{h \rightarrow 0} 2x + h \pause = 2x + 0 = 2x
  \end{align} \pause

  So we immediately get that: \pause
  \begin{enumerate}
    \item $f'(0) = 2 \cdot 0 = 0$ \pause
    \item $f'(1) = 2 \cdot 1 = 2$ \pause
    \item $f'(2) = 2 \cdot 2 = 4$
  \end{enumerate}
\end{frame}

\begin{frame}
  \frametitle{The Derivative Function III}

  \begin{alertblock}{Big Idea of Making it a Function}\index{derivative!as a function!why?}
    Here're two very important reasons why we might want to treat the derivative as a function and not just compute value by value: \pause
    \begin{enumerate}
      \item Speed. \pause Once we do the (slightly length) work of
        computing the general functional form of the derivative \pause
        all future derivatives at a point can be computed \textit{so quickly!} just by evaluation! \pause
      \item Information Extraction. \pause Without the equation of the function, it's \textit{impossible} to analytically study interesting properties of the derivative like: \pause
        \begin{enumerate}
          \item For which values is $f'(a)$ positive? (when is $f(x)$ increasing?) \pause
          \item For which values is $f'(a)$ zero? (important for maximizing/minimizing which we'll do in ~3/4 weeks) \pause
          \item For which values is $f'(a)$ undefined? (when can we not linearly approximate?) \pause
        \end{enumerate}
        but given an equation of a function it's fairly easy to answer all of those!
    \end{enumerate}
  \end{alertblock}
\end{frame}

\begin{frame}
  \frametitle{The Derivative Function IV}
  For example: at which $x$ values is the slope of the line tangent to $f(x) = x^3 - 2x + 1$ equal to zero? \pause

  Step 1: Translate to derivative question: \pause when is $f'(x) = 0$? \pause

  Step 2: Compute derivative: \pause
  \begin{align}
    f'(x) &= \pause \lim\limits_{h \rightarrow 0} \dfrac{f(x + h) - f(x)}{h} \pause \\
          &= \lim\limits_{h \rightarrow 0} \dfrac{(x + h)^3 - 2(x + h) + 1 - (x^3 - 2x + 1)}{h} \pause \\
          &= \lim\limits_{h \rightarrow 0} \dfrac{x^3 + 3x^2 h + 3xh^2 + h^3 - 2x - 2h + 1 - x^3 + 2x - 1}{h} \pause \\
          &= \lim\limits_{h \rightarrow 0} \dfrac{3x^2 h + 3xh^2 + h^3 - 2h}{h} \pause \\
          &= \lim\limits_{h \rightarrow 0} 3x^2 + 3xh + h^2 - 2 \pause \\
          &= 3x^2 + 3x \cdot 0 + 0^2 - 2 \pause = 3x^2 - 2
  \end{align}
\end{frame}

\begin{frame}
  \frametitle{The Derivative Function V}
  Last slide: $f'(x) = 3x^2 - 2$. \pause

  Step 3: Answer derivative question: \pause when is $f'(x) = 3x^2 - 2 = 0$? \pause

  \begin{gather}
    3x^2 - 2 = 0 \pause \\
    3x^2 = 2 \pause \\
    x^2 = \dfrac{2}{3} \pause \\
    x = \pm \sqrt{\dfrac{2}{3}}
  \end{gather} \pause

  Step 4: Translate derivative question answer to original question answer: \pause

  \begin{center}
    When $x$ is $\pm \sqrt{\dfrac{2}{3}}$, the slope of the line tangent to $f(x) = x^3 - 2x + 1$ is equal to zero.
  \end{center}
\end{frame}

\begin{frame}
  \frametitle{The Derivative Function VI}

  This is the general strategy that you can use to solve \textit{sooo} many questions about the behavior of functions! \pause

  We'll see that in the worksheet for today :)
\end{frame}

\begin{frame}
  \frametitle{The Derivative Function VII}

  Some conceptual questions: \pause
  \begin{enumerate}
    \item For some random function $f(x)$ whose \textit{domain} is all real numbers $\RR$ \pause must it also be the case that the domain of $f'(x)$ is also $\RR$? \pause
    \item How do the following conditions relate: \pause
      \begin{multicols}{2}
      \begin{enumerate}
        \item $f(x)$ is defined at $x = a$ 
        \item $\lim\limits_{x \rightarrow a} f(x)$ exists
        \item $f(a) = \lim\limits_{x \rightarrow a} f(x)$
        \item The domain of $f(x)$ is all real numbers $\RR$
        \item $\lim_{x \rightarrow a} \dfrac{f(x) - f(a)}{x - a}$ exists
        \item $f'(a)$ exists
        \item $f'(x)$ is defined at $x = a$
        \item $f(x)$ is differentiable at $x = a$
        \item $f'(x)$ is continuous at $x = a$
        \item $f'(x)$ is differentiable at $x = a$
        \item The domain of $f'(x)$ is all real numbers $\RR$
      \end{enumerate}
      \end{multicols} \pause
    \item Which of the above conditions relied on $f'(x)$ \textit{being a function} rather than being a limit defined at a point $x = a$?
  \end{enumerate}
\end{frame}

\begin{frame}
  \frametitle{Worksheet Time!}
  Let's get into groups and work on the worksheet!
\end{frame}

\begin{frame}
  \frametitle{End of Class!}
  That's all folks :)
\end{frame}
\end{document}
